\section{Conclusion}
\label{sec:conclusion}

We have presented EDEM, an agent-based framework that models supply
and demand as the realisation of a coupled stochastic process driven
by heterogeneous, error-prone agent valuations. The framework
recovers the classical Walrasian fixed point and a Miller-style
static premium mechanism as limiting cases, and admits a family of
out-of-equilibrium regimes -- band-stable, business-cycle, persistent
overshoot, persistent undershoot, runaway bubble, and constant
transition -- that are reachable from the same agent ruleset by
varying three parameters: the divergence-of-opinion scale
$\bar{\sigma}$, the seller patience timer, and the population
balancer coefficient $C_{b}$.

The eight-experiment Mesa replication of \cref{sec:experiments}
supports three substantive empirical claims. First,
\emph{textbook supply-and-demand schedules omit a state variable}:
seller patience and agent density both shift the realised market
price away from the equilibrium predicted by linear schedules, and
in opposite directions, while preserving the schedules themselves.
Second, \emph{bubbles do not require behavioural bias}; an unbiased
estimator combined with a max-bid clearing rule produces the same
multiplicative drift that observers attribute to investor optimism.
Third, \emph{linear balancer interventions cannot eliminate the
drift}: in the parameter range we explored, no setting of $C_{b}$
restores the textbook valuation, because the asymmetry is structural
to the bidding mechanism.

An immediate corollary is for valuation algorithms. A
machine-learning estimator trained on historical sale prices is
fitting an upper-order statistic of plausible sale prices, not their
mean; deployed at scale, it functions as a positively-biased
estimator that may provide the market with a coordinating signal
amplifying the very bubble it was meant to price. We do not advance
this as a complete account of any specific historical episode (see
\cref{sec:disc-ml}). The robust theoretical recommendation is that
clearing-price-anchored valuations be replaced with selection-aware
procedures (median rather than mean of plausible sale prices, or
censored-data likelihoods that account for the offer-acceptance
threshold) before being used as deployment-time point estimators.

EDEM's framework is intentionally general: the eight experiments
herein use a real-estate parameterisation, but the divergence of
opinion premium they explain has been documented across asset
classes, including initial public offerings
\citep{miller1977,doukas2006} and cryptocurrency markets. The Mesa
implementation is open-source and the eight experiments reproduce
to bit-equivalent figures from a single \texttt{pytest}-and-script
invocation; we hope this lowers the barrier to extending the
framework toward leveraged markets, adaptive agents, and
empirically-calibrated parameter sweeps that the present
exposition leaves for future work.
